\begin{figure}[tbp]
\centering
\begin{tikzpicture}[>=Stealth,font=\small,x=1cm,y=1cm,
    every node/.style={inner sep=1.5pt},
    model/.style={draw,rounded corners=3pt,fill=black!8},
    pt/.style={circle,fill,inner sep=0pt,minimum size=3.5pt},
    der/.style={->,decorate,decoration={snake,amplitude=0.6pt,segment length=5pt,post length=3pt}},
    dual/.style={->,dashed},
    opp/.style={densely dotted,thick},
    lb/.style={font=\scriptsize,inner sep=1pt},
    ttl/.style={font=\footnotesize,anchor=south}]
  \foreach \k in {0,1,2,3} {
    \begin{scope}[xshift=3.15*\k cm]
      \draw[black!50] (0,0) -- (2.95,0);
      \draw[black!25] (0,-1.25) rectangle (2.95,1.6);
    \end{scope}}
  \node[lb,anchor=east,align=right] at (-0.05,0.8) {visible\\ atoms};
  \node[lb,anchor=east,align=right] at (-0.05,-0.6) {hidden\\ atoms};
  % (a) (I0)
  \node[ttl] at (1.475,1.6) {(a) \CondI\ fails};
  \draw[model] (0.2,0.2) rectangle (1.6,1.4);
  \node[lb] at (0.55,1.2) {$\mu_\varphi$};
  \node[pt,label={[lb]below:fact $e$}] (ea) at (0.9,-0.7) {};
  \node[pt,label={[lb]right:$b$}] (ba) at (2.4,0.75) {};
  \draw[der] (ea) -- (ba);
  % (b) (U)
  \begin{scope}[xshift=3.15cm]
    \node[ttl] at (1.475,1.6) {(b) \CondU\ fails};
    \draw[model] (0.2,0.2) rectangle (1.6,1.4);
    \node[lb] at (0.6,1.2) {$\Cl_\varphi(a)$};
    \node[pt,label={[lb]left:$a$}] (ab) at (0.9,0.65) {};
    \node[draw,double,rounded corners=1pt,lb,inner sep=2pt,align=center] (fb) at (2.2,-0.7) {forbidden\\ root};
    \draw[der] (ab) -- (fb);
  \end{scope}
  % (c) (II)
  \begin{scope}[xshift=6.3cm]
    \node[ttl] at (1.475,1.6) {(c) \CondII\ fails};
    \draw[model] (0.2,0.2) rectangle (1.6,1.4);
    \node[lb] at (0.6,1.2) {$\Cl_\varphi(a)$};
    \node[pt,label={[lb]left:$a$}] (ac) at (0.9,0.65) {};
    \node[pt,label={[lb]right:$b$}] (bc) at (2.4,0.75) {};
    \draw[der] (ac) .. controls (1.6,-0.9) and (2.2,-0.6) .. (bc);
  \end{scope}
  % (d) collision
  \begin{scope}[xshift=9.45cm]
    \node[ttl] at (1.475,1.6) {(d) two seeds collide};
    \draw[model] (0.15,0.2) rectangle (1.75,1.4);
    \node[lb] at (0.45,1.2) {$S$};
    \node[pt,label={[lb]above:$a$}] (ad) at (0.5,0.6) {};
    \node[pt,label={[lb]above:$(q,v)$}] (qd) at (1.35,0.85) {};
    \node[pt,label={[lb]below:$b$}] (bd) at (1.0,-0.65) {};
    \node[pt,label={[lb]below:$\bar b$}] (nbd) at (2.0,-0.65) {};
    \node[pt,label={[lb]above:$(\bar q,v\delta)$}] (rd) at (2.4,0.95) {};
    \draw[der] (ad) -- (bd);
    \draw[der] (qd) -- (nbd);
    \draw[opp] (bd) -- (nbd);
    \draw[dual] (bd) .. controls (1.5,0.0) and (2.1,0.4) .. (rd);
  \end{scope}
\end{tikzpicture}
\caption{When both systems are satisfiable, the four ways in which a model $S$
of $\varphi$ can fail to extend to a model of $\psi$, in the proof of
\cref{thm:ch-characterization}. Upper halves
hold visible atoms, lower halves hidden ones; wavy arrows are derivations of
$\psi$, and the shaded box is a model of $\varphi$. (a) A fact of $\psi$
derives a visible atom outside the least model $\mu_\varphi$. (b) A seed $a$
derives a forbidden root; then the least model $\Cl_\varphi(a)$ of $\varphi$
that contains $a$ cannot be extended. (c) A seed $a$ derives a visible atom
outside $\Cl_\varphi(a)$. (d) Two seeds of $S$ derive opposite atoms $b$ and
$\bar b$, possibly hidden. Collision return (\cref{lem:lf-collision}) reverses
the second derivation by duality (dashed): the single seed $a$ derives the
opposite of the extension $(q,v\delta)\in S$ of the second seed, an atom
outside $S\supseteq\Cl_\varphi(a)$, so this is case~(c). If $a$ is a fact of
$\psi$ instead, it is case~(a).}
\Description{Four panels, each split by a horizontal line into visible atoms
above and hidden atoms below. In each panel a shaded box is a model of phi.
Panel a: a fact e of psi, below the line, derives by a wavy arrow a visible
atom b outside the box labeled mu phi. Panel b: a seed a inside the box,
labeled closure of a, derives a doubly framed forbidden root below the line.
Panel c: a seed a inside the box derives, through hidden atoms, a visible atom
b outside the box. Panel d: two seeds a and q v inside the box labeled S
derive the hidden atoms b and b bar, joined by a dotted line; a dashed arrow from b
leads to the visible atom q bar, v delta outside the box.}
\label{fig:ch-conditions}
\end{figure}
